\subsection{Cosmologia da delegação}\label{sec:orbits}
% ============================================================

A propagação $P$ (Seção~\ref{sec:algebra}) é um \emph{fluxo}: um \emph{token}
--- uma capacidade com profundidade restante --- viaja pelo grafo, perdendo
uma unidade de $\dep$ a cada salto. Como o grafo pode ter ciclos
(Seção~\ref{sec:dag}), uma trajetória pode reencontrar um nó já visitado: é
uma \textbf{órbita}. A pergunta dinâmica é se essas órbitas terminam.

\paragraph{$\dep$ é uma energia.} A profundidade restante é uma função de
Lyapunov \emph{estrita}: cada aresta válida a decrementa em exatamente $1$
(Definição~\ref{def:att}). Logo, ao percorrer um ciclo de comprimento $\ell$,
o token volta ao nó de partida com $\dep$ menor em $\ell$ --- a órbita não se
fecha, \textbf{espirala para dentro} até $\dep=0$, onde o fluxo para
(Lema~\ref{lem:energy}). A aciclicidade é dispensável; a terminação vem da
\emph{dissipação} de $\dep$.

\paragraph{A analogia de mecânica celeste.} Esta é a versão \emph{dissipativa}
de um sistema conservativo bem conhecido. Tomando o espaço de álgebras
hipercomplexas com a métrica de de Sitter 2D
$ds^2=-\cos^2\psi\,d\theta^2+d\psi^2$,\footnote{Veja
\texttt{hiper/circular/paper\_metrica\_quantica\_algebras.tex}: a métrica
quântica como geometria das álgebras. A seção euclidiana ($\theta\to-i\theta_E$)
é a esfera $S^2$, onde as geodésicas são círculos máximos --- órbitas
fechadas.} a energia de Noether $E=\cos^2\psi\,\dot\theta$ (do vetor de Killing
$\theta$) é \textbf{conservada}, e a geodésica é uma \emph{órbita fechada}. A
delegação tem a mesma forma com um termo a mais: a atenuação é o
\textbf{atrito} que dissipa a energia. Onde a órbita de de Sitter se fecha
(energia constante), a órbita de delegação decai (Figura~\ref{fig:orbits}).
Coerentemente, $\dep=\infty$ --- reservado às autoridades intrínsecas --- é o
caso \emph{sem atrito}: energia conservada, órbita que nunca decai.

\begin{remark}[Um universo de vida curta]\label{rem:vida}
A delegação é, então, \emph{o universo de de Sitter num meio dissipativo}: a
mesma geometria, com atrito. O universo conservativo ($\dep=\infty$) é eterno;
toda concessão finita é um universo de \textbf{vida curta}. E há uma sutileza
que aguça a imagem: o atrito da delegação é \emph{a taxa constante}
($-1$/salto), não viscoso ($\propto$ velocidade). O arrasto viscoso dá decaimento
\emph{exponencial} --- assintótico, nunca para de fato; o atrito de taxa
constante (atrito \emph{seco}, de Coulomb) dá um \textbf{tempo de vida finito e
exato}: $\tau_{\mathrm{vida}}=\dep$ saltos. A órbita não desvanece --- ela
\emph{para}, de forma abrupta, após exatamente $\dep$ saltos. É o que torna a
não-escalada não apenas verdadeira no limite, mas \emph{decidível em tempo
finito}.
\end{remark}

\begin{proposition}[Não há órbita infinita]\label{prop:orbit}
Em qualquer grafo de delegação \emph{finito} (com ciclos), e para qualquer
capacidade emitida com $\dep<\infty$: toda trajetória de propagação tem
comprimento $\le\dep$; o alcance estabiliza no menor ponto fixo
$\mu P=\bigsqcup_{k\le\dep}P^k(\{\mathrm{orig}\})$; e a iteração de Kleene
converge em $\le\dep$ passos.
\end{proposition}
\begin{proof}
$\dep$ decresce $1$ por salto (variante bem-fundado), logo nenhuma cadeia
excede $\dep$; como $P$ é monótona num reticulado finito e a difusão satura
quando $\dep$ se esgota, $P^{\dep}=P^{\dep+1}=\mu P$.
\end{proof}

\paragraph{Verificação a escala.} O experimento
\texttt{doc/orbits/delegation\_orbits.py} constrói um grafo de delegação
\emph{grande e com ciclos} ($3000$ nós, $12{,}012$ arestas, $2880$ deles num
componente fortemente conexo --- quase tudo em ciclos) e confirma a
Proposição~\ref{prop:orbit} numericamente, com $\dep=20$:
\begin{itemize}[nosep,leftmargin=1.3em]
  \item \textbf{terminação apesar dos ciclos}: a maior cadeia de re-delegação
        realizada é $9\le\dep$;
  \item \textbf{energia bem-fundada}: toda transição decrementa $\dep$ em $1$;
  \item \textbf{ponto fixo $\mu P$}: o alcance ($2938$ nós) satura em $k=10$;
  \item \textbf{monotonicidade}: o alcance é não-decrescente em $k$.
\end{itemize}

  \begin{figure}[h]
  \centering
  \IfFileExists{orbits/delegation_orbits.png}{%
    \includegraphics[width=\textwidth]{orbits/delegation_orbits.png}%
  }{%
    \begin{tikzpicture}[
        every node/.style={font=\small,align=center},
        orb/.style={draw=ourovivo,fill=creme,rounded corners=2pt,minimum width=3.2cm,minimum height=1.1cm},
      ]
      \node[orb] (a) {órbita\\dissipativa ($\dep\downarrow$)};
      \node[orb, right=1.2cm of a] (b) {Kleene\\$\mu P$};
      \node[orb, right=1.2cm of b] (c) {fail-closed\\$\varnothing$};
      \draw[->,ourovivo,thick] (a)--(b); \draw[->,ourovivo,thick] (b)--(c);
    \end{tikzpicture}%
  }
  \caption{Órbitas de delegação \emph{vs} geodésica conservativa. Com $\dep<\infty$ a
  propagação termina; sem figura externa, esquema tikz do fluxo dissipativo.}
\label{fig:orbits}
\end{figure}

\begin{remark}[Não é identificação física, é analogia dinâmica]
O caso conservativo (de Sitter/esfera) e o dissipativo (delegação) diferem
\emph{exatamente} pelo termo de atrito $=$ atenuação. É essa dissipação que
torna o sistema seguro a escala --- com ciclos e tudo, sem precisar de DAG.
\end{remark}

\begin{remark}[Duas leituras: matemática e observabilidade]\label{rem:leituras}
Esta seção admite duas leituras. \textbf{Leitura A} (matemática): o sistema
lembra uma dinâmica dissipativa em de Sitter. \textbf{Leitura B} (produto): o
estado de autorização do \SO\ forma um \emph{universo observável} --- no lab
\texttt{conversa/}, via \texttt{observar.py}, \texttt{mede\_flocos.py} e
\texttt{cadeias.py --mede}. Quem cura dialogos não lê provas de não-escalada,
mas vê instantaneamente quando a autoridade \emph{nasce num nó curado},
\emph{flui} pelo grafo de turnos, \emph{perde energia} ($\dep\downarrow$),
\emph{forma órbitas} (ciclos Koch inofensivos) e \emph{para} no horizonte
$\dep=0$. Grandezas úteis: overlap médio, frames $\Phi$, atratores únicos,
taxa de acerto por cadeia. Anomalia --- colisão persistente, queda em
\texttt{Oi oi oi!} --- é assinatura geométrica, não conteúdo lexical.
\end{remark}

% ============================================================
\subsection{Nível de campo: a delegação em de Sitter 4D}\label{sec:field}
% ============================================================

Até aqui a condição $c$ seleciona \emph{linhas}. Mas o CASL também restringe
\emph{campos} (colunas): uma regra pode conceder ``ler \texttt{Customer.name,
\,Customer.email}'' sem dar \texttt{Customer.document}. Isto é uma segunda
dimensão.

\begin{definition}[Capacidade a nível de campo]
Seja $\F$ o conjunto de campos (de um subject). Uma capacidade ganha um
\emph{conjunto de campos} $\varphi\subseteq\F$:
$\kappa=(a,s,c,\varphi,\iota,\dep)$, com extent
$\ext{\kappa}=\ext{c}\times\varphi\subseteq\R\times\F$. A autorização efetiva
vive agora em $2^{\R\times\F}$.
\end{definition}

\paragraph{O \emph{lift} é de graça (a álgebra é agnóstica ao portador).} Tudo
o que provamos --- deny-overrides (Prop.~\ref{prop:deny}), atenuação e
não-escalada (Teor.~\ref{thm:noesc}), $\dep$ e o ponto fixo $\mu P$, a
cosmologia (Seção~\ref{sec:orbits}) --- usou \emph{apenas} a estrutura de
dioid booleano de $2^{X}$, \textbf{nunca} o portador $X$ específico. Trocar
$X=\R$ por $X=\R\times\F$ não muda uma linha das demonstrações. A atenuação
agora estreita \emph{dois} eixos independentes: $\ext{c'}\subseteq\ext{c}$
(linhas) \emph{e} $\varphi'\subseteq\varphi$ (campos). O \emph{pushdown}
projeta a condição no \texttt{where} (linhas) e o $\varphi$ no \texttt{select}
(campos) --- o que a extensão RLS já faz com \texttt{fields}.

\paragraph{A geometria sobe duas dimensões.} O espaço dinâmico do nível de
linha era $\dS_2$ (Seção~\ref{sec:orbits}): a fase $t$ do fluxo de delegação
mais a energia $\dep$ (o $\rho$). Os campos acrescentam suas \emph{direções}
--- uma esfera $S^2$ de configurações de campo $(\alpha,\beta)$, à la as
esferas de direções de álgebra do paper. A forma estática de de Sitter
generaliza exatamente assim:
\begin{equation}
  \dS_n:\quad ds^2=-\cos^2\!\rho\,dt^2+d\rho^2+\sin^2\!\rho\,d\Omega^2_{n-2},
  \label{eq:dSn}
\end{equation}
onde $(t,\rho)=(\text{fluxo, energia }\dep)$ e $d\Omega^2_{n-2}$ são as
direções de campo. O nível de linha é \eqref{eq:dSn} com $n=2$ (a métrica do
paper, $d\Omega_0$ trivial); o nível de campo com uma esfera de campo $S^2$ é
$n=4$:
\[
  ds^2=-\cos^2\!\rho\,dt^2+d\rho^2+\sin^2\!\rho\,(d\alpha^2+\sin^2\!\alpha\,d\beta^2).
\]

\begin{theorem}[de Sitter 4D suporta o nível de campo]\label{thm:dS4}
A métrica \eqref{eq:dSn} com $n=4$ é um espaço de Einstein de curvatura
constante --- de Sitter $\dS_4$ de raio $\ell=1$:
\[
  R_{\mu\nu}=(n-1)\,g_{\mu\nu}=3\,g_{\mu\nu},\qquad R=n(n-1)=12.
\]
\end{theorem}
\begin{proof}[Verificação]
Cálculo do tensor de Ricci em \texttt{doc/orbits/desitter4d.py} (Christoffel
$\to$ Riemann $\to$ Ricci), conferido numericamente em $12$ pontos (resíduo
$\sim10^{-16}$), \emph{exatamente} a metodologia do paper. O caso $n=2$
reproduz $\dS_2$ ($R_{\mu\nu}=g_{\mu\nu}$, $R=2$) como controle. Em geral,
$k$ dimensões de campo dão $\dS_{2+k}$.
\end{proof}

\begin{remark}[A cosmologia se completa]
O horizonte $\rho=\pi/2$ (energia $\dep$ esgotada, o quaternion do paper)
continua sendo o fim de vida da órbita (Seção~\ref{sec:orbits}); a esfera de
campo $\sin^2\!\rho\,d\Omega^2_2$ acopla as direções de campo à mesma energia.
A autorização a nível de campo é, então, um universo $\dS_4$ em meio
dissipativo --- a mesma física do nível de linha, com duas direções
espaciais a mais. A escolha do nível (linha vs campo) é só a dimensão do
cosmos; a segurança (atenuação, não-escalada, vida finita) é idêntica em
qualquer $n$.
\end{remark}

% ============================================================
